% ======================================================================
% Appendix -- Notation and glossary
% Accessibility requirement (researcher-set): every recurring term and
% symbol, its plain-English meaning, and where it is first used.
% ======================================================================
\section{Notation and Glossary}
\label{app:glossary}

\Cref{tab:glossary} collects the paper's recurring terms and symbols.
Each entry gives the plain-English meaning used throughout, anchored in
the molecule running example (atoms $=$ nodes $=$ rank 0, bonds $=$
edges $=$ rank 1, rings $=$ faces $=$ rank 2), and the section where
the term is first defined.

\begin{table}[h]
\centering
\caption{Glossary of terms and symbols.}
\label{tab:glossary}
\small
\begin{tabular}{p{0.20\linewidth}p{0.56\linewidth}l}
\toprule
Term / symbol & Plain meaning & First use \\
\midrule
cell & A node or a group of nodes stored as one object: an atom, a
  bond, a ring. & \cref{sec:prelim-cc} \\
rank, $\rk$ & Which level a cell lives on: atoms 0, bonds 1, rings 2.
  & \cref{sec:prelim-cc} \\
combinatorial complex, $\mathcal{C}$ & The data structure holding all
  cells of all ranks for one input. & \cref{sec:prelim-cc} \\
lifting & Converting a plain graph into a complex, \eg promoting a
  molecule's chordless cycles to ring cells. & \cref{sec:prelim-cc} \\
neighborhood, $N$ & A message-passing route: who talks to whom (\eg
  atom$\to$bond, bond$\to$ring). & \cref{sec:prelim-gccn} \\
GCCN & Generalized combinatorial complex network: a TDL model built
  from one submodule per neighborhood. & \cref{sec:prelim-gccn} \\
neighborhood-indexed component & The architecture piece a model has
  one of per neighborhood: a GCCN route, a HOPSE encoding block.
  & \cref{sec:framework-game2} \\
game, $v$ & An assignment of a score to every subset of players.
  & \cref{sec:prelim-shapley} \\
coalition, $S$ & A subset of players (cells, or components).
  & \cref{sec:prelim-shapley} \\
marginal contribution & How much the score changes when one player
  joins a coalition. & \cref{sec:prelim-shapley} \\
Shapley value, $\phi_i$ & A player's average marginal contribution over
  all coalitions; the unique fair credit split.
  & \cref{sec:prelim-shapley} \\
masking & Neutralizing the players outside a coalition so the model
  behaves as if they were absent. & \cref{sec:framework-game1} \\
exact (retraining) game & Score of a coalition $=$ validation accuracy
  after training from scratch with exactly those components.
  & \cref{sec:framework-game2} \\
cheap (masked-validation) game, $v_{\mathrm{prune}}$ & Score of a
  coalition $=$ the brief run's validation accuracy with the other
  components removed; one forward pass per coalition.
  & \cref{sec:framework-game2} \\
brief run & The all-components model trained briefly
  (\unfrozen{\ResStemEpochs} epochs) before attribution.
  & \cref{sec:framework-game2} \\
candidate sequence & The nested sequence of shrinking architectures
  produced by backward elimination; one candidate per size.
  & \cref{sec:method-pipeline} \\
candidate & One architecture in the candidate sequence.
  & \cref{sec:method-pipeline} \\
stop candidate & The smallest candidate whose next removal would cost
  more than the noise floor. & \cref{app:selection-ladder} \\
budget $B$ & How many candidates are trained to completion before
  validation picks; the practitioner's dial.
  & \cref{sec:method-pipeline} \\
binomial floor & The validation set's own noise level,
  $\sqrt{p(1-p)/m}$; the elimination's stop threshold.
  & \cref{sec:method-pipeline} \\
substitute-commitment & The brief run's arbitrary choice between two
  components that can do the same job; attribution reports the choice,
  not the interchangeability. & \cref{app:selection-commitment} \\
menu & The \unfrozen{\ResMenuEntries} published fixed architectures
  from the literature, treated as one baseline at its true cost.
  & \cref{sec:method-pipeline} \\
sweep & Training one model per candidate configuration and keeping the
  best; the procedure the pipeline replaces. & \cref{sec:prelim-gccn} \\
\bottomrule
\end{tabular}
\end{table}
